% ---- pg-378 （书页 366） ----

\[
\max_{\mu_i,\sigma_{ii}}(2\pi)^{-n/2}\sigma_{ii}^{-n/2}\exp\left[-\sum_{j=1}^{n}(x_{ji}-\mu_i)^2/2\sigma_{ii}\right],\quad \text{所以, } \hat{\mu}_i=n^{-1}\sum_{j=1}^{n}x_{ji}
\]
且 $\hat{\sigma}_{ii}=(1/n)\sum_{j=1}^{n}(x_{ji}-\bar{x}_i)^2$. 除数 $n$ 在 $\Lambda$ 中抵消, 故可用 $\mathbf{S}$.

(b) 在 $H_0$ 下验证 $\hat{\sigma}^2=\left[\sum_{j=1}^{n}(x_{j1}-\bar{x}_1)^2+\cdots+\sum_{j=1}^{n}(x_{jp}-\bar{x}_p)^2\right]/np$.

在统计量中去掉除数 $n$, 再验证仍然可用 $\mathbf{S}$. 利用结论 5.2 计算 $\hat{\chi}^2$ 的自由度.

\mvbold{以下练习需用计算机.}

\noindent 8.10\quad \hyperref[tab:8.4]{表 8.4} 给出在纽约股票交易所交易的五支股票的周回报率(见网页 www.prenhall.com/statistics 上的股票价格数据).

(a) 构造样本协方差矩阵 $\mathbf{S}$, 并求\hyperref[eq:8-20]{式(8-20)}的样本主成分. (注意, 样本均值向量 $\bar{\mathbf{x}}$ 已在例 8.5 中列出了.)

(b) 确定样本总方差中用前三个主成分解释的比例. 解释这些成分.

(c) 前三个总体成分 $Y_1,Y_2$ 和 $Y_3$ 的方差为 $\lambda_1,\lambda_2$ 和 $\lambda_3$, 构造它们的庞弗罗尼 $90\%$ 联合置信区间.

(d) 已知 (a)$\sim$(c) 的结果, 你认为要综合股票回报率数据, 在小于五维空间中是否可行? 试解释.

\begin{mvtable}
\phantomsection\label{tab:8.4}\mvtabcap{表 8.4\quad 股票价格数据(周回报率)}
\begin{longtable}{lrrrrr}\toprule
周 & JP Morgan & Citibank & Wells Fargo & Royal Dutch Shell & Exxon Mobil\\
\midrule
\endfirsthead
\multicolumn{6}{c}{\zihao{6}续表}\\
\toprule
周 & JP Morgan & Citibank & Wells Fargo & Royal Dutch Shell & Exxon Mobil\\
\midrule
\endhead
\bottomrule
\endlastfoot

1 & 0.013 03 & $-0.007$ 84 & $-0.003$ 19 & $-0.044$ 77 & 0.005 22\\
2 & 0.008 49 & 0.016 69 & $-0.006$ 21 & 0.01196 & 0.013 49\\
3 & $-0.017$ 92 & $-0.008$ 64 & 0.010 04 & 0 & $-0.006$ 14\\
4 & 0.021 56 & $-0.003$ 49 & 0.017 44 & $-0.028$ 59 & $-0.006$ 95\\
5 & 0.010 82 & 0.003 72 & $-0.010$ 13 & 0.029 19 & 0.040 98\\
6 & 0.010 17 & $-0.012$ 20 & $-0.008$ 38 & 0.013 71 & 0.002 99\\
7 & 0.011 13 & 0.028 00 & 0.008 07 & 0.030 54 & 0.003 23\\
8 & 0.048 48 & $-0.005$ 15 & 0.018 25 & 0.006 33 & 0.007 68\\
9 & $-0.034$ 49 & $-0.013$ 80 & $-0.008$ 05 & $-0.029$ 90 & $-0.010$ 81\\
10 & $-0.004$ 66 & 0.020 99 & $-0.006$ 08 & $-0.020$ 39 & $-0.012$ 67\\
$\vdots$ & $\vdots$ & $\vdots$ & $\vdots$ & $\vdots$ & $\vdots$\\
94 & 0.037 32 & 0.035 93 & 0.025 28 & 0.058 19 & 0.016 97\\
95 & 0.023 80 & 0.003 11 & $-0.006$ 88 & 0.012 25 & 0.028 17\\
96 & 0.025 68 & 0.052 53 & 0.040 70 & $-0.031$ 66 & $-0.018$ 85\\
97 & $-0.066$ 06 & 0.008 63 & 0.005 84 & 0.044 56 & 0.030 59\\
98 & 0.021 74 & 0.022 96 & 0.029 20 & 0.008 44 & 0.031 93\\
99 & 0.003 37 & $-0.015$ 31 & $-0.023$ 82 & $-0.001$ 67 & $-0.017$ 23\\
100 & 0.003 36 & 0.002 90 & $-0.003$ 05 & $-0.001$ 22 & $-0.009$ 70\\
101 & 0.017 01 & 0.009 51 & 0.018 20 & $-0.016$ 18 & $-0.007$ 56\\
102 & 0.010 39 & $-0.002$ 66 & 0.044 43 & 0.002 48 & $-0.016$ 45\\
103 & $-0.012$ 79 & $-0.014$ 37 & $-0.018$ 74 & $-0.004$ 98 & $-0.016$ 37\\
\end{longtable}
\end{mvtable}
